\begin{abstract}
Long-horizon coding agents are usually made stronger by giving them more context, as if the binding constraint were how much of a repository fits in the window.
We treat the prompt as a compiled artifact instead: typed sections under per-kind share budgets, a locator card that names where to look and which test failed, a ledger of what the model was actually shown, and mechanical gates that decide when the agent must stop revising and submit.
On a 16-instance\fact{A1} subset of SWE-bench Verified with DeepSeek V4.1 Flash, evidence delivery changed the resolve rate from 5/16\fact{A3} to 10/16\fact{A2} in a paired run where context assembly was the only variable.
On a single long task, six process constraints in the system prompt left MiniMax-M3.1-Flash at 0/3\fact{E5}; once the locator card was actually rendered and a mechanical cap allowed two patch versions, the same task closed at 5/5\fact{E8}.
Replaying those six tasks' recorded requests, our arm spent 73,051\fact{D2} output tokens per task against 12,495\fact{D2} for the official mcode CLI.
The scope is narrow: one model family per experiment, small task sets, one run per arm.
\end{abstract}

% FACTS-USED: A1, A2, A3, D2, E5, E8
